% Source: https://artofproblemsolving.com/community/c3h2382296_equation_that_hits_3_different_points
% Stable identity: aops-topic-2382296-q01
\textbf{(a)} The right side is a polynomial of degree at most $1$ in $x$, so its graph is a (non-vertical) line. At $x=x_1$ it equals $y_1$, and at $x=x_2$ it equals $y_1+(y_2-y_1)=y_2$. So it is the line through the two points.

\textbf{(b)} Build $q$ from three pieces, each vanishing at two of the given $x$-values (Lagrange's form):
\[q(x)=y_1\dfrac{(x-x_2)(x-x_3)}{(x_1-x_2)(x_1-x_3)}+y_2\dfrac{(x-x_1)(x-x_3)}{(x_2-x_1)(x_2-x_3)}+y_3\dfrac{(x-x_1)(x-x_2)}{(x_3-x_1)(x_3-x_2)}.\]
The denominators are nonzero because the $x_i$ are distinct. Each term has degree at most $2$, so $q$ does too. At $x=x_1$ the second and third terms vanish and the first equals $y_1$; similarly $q(x_2)=y_2$ and $q(x_3)=y_3$.

\textbf{Uniqueness.} If $q_1$ and $q_2$ both work, then $d=q_1-q_2$ has degree at most $2$ and the three roots $x_1,x_2,x_3$. A nonzero polynomial of degree at most $2$ has at most $2$ roots (by the factor theorem each root $x_i$ gives a factor $x-x_i$), so $d=0$ and $q_1=q_2$.

\textbf{(c)} The coefficient of $x^2$ in $q$ is
\[\dfrac{y_1}{(x_1-x_2)(x_1-x_3)}+\dfrac{y_2}{(x_2-x_1)(x_2-x_3)}+\dfrac{y_3}{(x_3-x_1)(x_3-x_2)}.\]
It is $0$ exactly when the three points are collinear: if they lie on a line $y=mx+c$, that line is a polynomial of degree at most $1$ through them, so by uniqueness it is $q$; conversely if $\deg q\le1$ its graph is a line through the three points. So $q$ is a genuine quadratic exactly when the three points are \emph{not} collinear.

(For example, through $(0,1)$, $(1,3)$, $(2,7)$ the formula gives $q(x)=x^2+x+1$.)
